\section{Simulation Protocols and Metrics}
\label{app:protocols}

\paragraph{Scope.}
This appendix provides supporting technical detail for the corresponding main-text sections. Essential claims and results remain in the main paper.

\subsection{Simulation conditions}

All reported experiments use one focal active-inference agent with four partners ($K=4$), per-partner generative models, and environment-side partner policies (Appendix~\ref{app:pomdp}). Conditions differ by payoff regime, assignment mode, and scheduled partner changes.

\textbf{Binary random-assignment baseline.} In the binary trust game (Table~\ref{tab:payoff} in the main text), partners are assigned randomly each round and the agent chooses cooperate or defect. Stochastic partner-type drift uses $p_{\mathrm{switch}}=0.05$ unless disabled for scheduled-switch scenarios.

\textbf{Open graded decision regime.} The graded regime replaces the binary action with a discretized investment level while keeping the same latent type--stance structure. This is the primary deployment-dissociation setting (Section~\ref{sec:deployment}): affect, no-affect, and tracked-only variants are compared on policy entropy and cumulative payoff.

\textbf{Partner-choice regime.} The agent first selects a partner, then selects an action toward that partner each round. Partner selection is where engagement reorganization and social-allocation readouts are defined (Section~\ref{sec:partner_choice}).

\textbf{Focal-switch locality probe.} A partner-choice graded environment with type volatility disabled ($p_{\mathrm{switch}}=0$) and a single scheduled stance switch on partner P0 tests whether local versus shared $\beta$ tracking preserves the surprise-over-reward dissociation (Section~\ref{sec:model_fitness}).

\textbf{Abrupt betrayal.} A partner-choice graded environment with scheduled stance betrayal on a previously cooperative partner tests post-switch reallocation, entropy, payoff, and joint accuracy (Section~\ref{sec:betrayal}).

\textbf{Precision-dynamics perturbations.} Clinical-like computational variants (alexithymia-like, borderline-like, depression-like) perturb $\beta_k$ amplitude and stability in the graded decision regime without claiming clinical validity (Appendix~\ref{app:extended_results}).

\textbf{Phenotype program (Exp A--D).} Parameter sweeps and factorial designs vary precision gain $\alpha$ and prior over model fitness to define computational trust-calibration profiles in open graded, betrayal, partner-choice, forgiveness, and mixed-volatility environments (Section~\ref{sec:phenotypes}).

\subsection{Seed counts and episode lengths}

Unless otherwise stated in the main text, episode length is 200 rounds. The abrupt-betrayal condition uses 120 rounds because post-switch readouts dominate interpretation. Phenotype and mixed-volatility tables use 20 seeds per condition. Central behavioral confirmations target 30 seeds per condition.

Exploratory diagnostic runs for the H0--H6 spine used three seeds per variant for exploratory validation of the surprisal-based precision tracker and centered cross-partner policy selection; those numbers remain underpowered and are not treated as confirmation-scale evidence. Current manuscript headline values drawn from exploratory diagnostic runs are flagged for replacement where 30-seed confirmation is pending (Section~\ref{sec:results}).

Representative configured scales:

\begin{table}[H]
\centering
\caption{Representative episode lengths and replication counts by experiment family.}
\label{tab:protocol_scales}
\begin{tabular}{llrr}
\toprule
Experiment family & Setting & Rounds & Seeds \\
\midrule
H0/H2 open graded deployment & graded, random or open choice & 200 & 30 (confirm.) / 3 (explor. diag.) \\
H1 model fitness & graded / reliability probes & 200 & 30 (confirm.) / 3 (explor. diag.) \\
H3 focal-switch locality & partner choice, single switch & 100 & 30 (confirm.) / 3 (explor. diag.) \\
H4 partner choice & partner choice & 200 & 30 (confirm.) / 3 (explor. diag.) \\
H5 abrupt betrayal & partner choice, scheduled betrayal & 120 & 30 (confirm.) / 3 (explor. diag.) \\
H6 clinical dynamics & graded perturbations & 200 & 30 (confirm.) / 3 (explor. diag.) \\
Exp A--D phenotypes & sweeps / factorial / forgiveness / mixed volatility & 200 & 20 \\
\bottomrule
\end{tabular}
\end{table}

\subsection{Betrayal schedule}

The canonical abrupt-betrayal protocol (H5 \texttt{betrayal\_choice}) uses graded payoffs, agent-chosen partners, four partners with fixed initial types and stances, and no stochastic type switching ($p_{\mathrm{switch}}=0$). Partner P0 begins as a cooperator in a trusting stance. At round 31, P0 undergoes a scheduled stance switch to hostile while other partners retain their initial configurations. Primary readouts are post-switch cumulative payoff, policy entropy, partner reallocation, and joint partner-type accuracy over the remaining episode.

The phenotype betrayal environment used in Exp A and Exp B applies the same single-switch structure within 200-round episodes for recovery-time and quadrant metrics. Exp C (forgiveness) extends this schedule: rounds 1--80 cooperative baseline, rounds 81--120 betrayal on P0, rounds 121--200 reversion to cooperative/trusting on P0.

\subsection{Mixed-volatility schedule}

The mixed-volatility environment (Exp D; Section~\ref{sec:phenotypes_mixed}) places four partners with heterogeneous volatility in the same partner-choice episode over 200 rounds:

\begin{itemize}
\item \textbf{P0:} stationary cooperator (no scheduled change).
\item \textbf{P1:} stationary exploiter (no scheduled change).
\item \textbf{P2:} abrupt shift---cooperative/trusting pre-switch, hostile/exploiter from the scheduled switch onward (round 101 in the implemented scenario).
\item \textbf{P3:} gradual drift---stance and type transitions spread across mid-episode rounds (neutral stance from round 51; hostile/exploiter shift from round 151 in the implemented scenario).
\end{itemize}

Stochastic type switching is disabled ($p_{\mathrm{switch}}=0$). Conditions compare default, low-$\alpha$, high-$\alpha$, and no-affect variants at 20 seeds each.

\subsection{Partner-choice allocation metrics}

Partner-choice readouts quantify how precision reshapes \emph{who} the agent engages, not only how it acts toward a fixed partner.

\textbf{Selection share.} Per-partner selection rate is the fraction of rounds on which each partner is chosen. Pending confirmation-scale reporting includes affect versus no-affect selection shares by partner (Section~\ref{sec:partner_choice}).

\textbf{Selection concentration.} The Gini coefficient of the partner-selection distribution summarizes how concentrated engagement is across the social portfolio (used in phenotype Exp A and Exp B).

\textbf{Post-betrayal allocation.} Post-switch windows report reallocation away from the switched partner and toward alternatives, including post-betrayal P0 selection and high-investment rates in phenotype analyses.

\textbf{Policy entropy under choice.} Mean $q_\pi$ entropy in partner-choice regimes captures how sharply the agent commits to partner--action policies (reported comparatively for affect versus no-affect controls).

\subsection{Entropy, payoff, and accuracy readouts}

\textbf{Total payoff.} Cumulative focal-agent payoff over the episode (or over defined post-switch windows for betrayal analyses).

\textbf{Policy entropy.} Mean $q_\pi$ entropy (`q\_pi\_entropy` in logged outputs) measures action-commitment sharpness; lower entropy indicates sharper policy selection under precision modulation.

\textbf{Joint partner-type accuracy.} Mean joint correctness of inferred partner type (`mean\_joint\_accuracy`) summarizes epistemic state inference; it is reported separately from payoff because the betrayal result can show higher payoff with lower accuracy when reallocation succeeds.

\textbf{$\beta_k$ range.} Per-partner or episode-aggregated range of $\beta_k$ support summarizes precision-dynamics amplitude in phenotype and perturbation tables.

\textbf{Model-fitness dissociation.} Partner-local correlations between mean precision and action surprisal versus realized payoff test whether $\beta_k$ tracks predictive reliability rather than partner value (Section~\ref{sec:model_fitness}).

\textbf{Mixed-volatility discrimination metrics.} The discrimination index is the Pearson correlation between per-partner $\beta_k$ trajectory and per-partner behavioral stability (higher indicates better separation of stable versus shifting partners). The P0 false-positive rate is the rolling rate at which engagement with the stationary cooperator P0 falls more than 15\% below its early baseline despite P0 never changing.

\textbf{Betrayal timing metrics.} Recovery time and related latency readouts measure rounds after a switch until selection or payoff returns toward baseline thresholds (phenotype Exp A/B and betrayal analyses).

\subsection{Placeholder results awaiting final runs}

The following readouts remain pending confirmation-scale values. Main-text placeholders and TODO comments are retained in Section~\ref{sec:results}; this list records what still awaits final runs without filling numbers.

\begin{itemize}
\item \textbf{Shared-$\beta$ ablation} (Section~\ref{sec:shared_beta}): partner-local versus shared precision on payoff, entropy, reliability-specificity, partner allocation, and false-positive suppression in mixed-volatility or focal-switch regimes.
\item \textbf{Partner-choice allocation rates} (Section~\ref{sec:partner_choice}): selection share by partner under affect and no-affect, and whether aggregate payoff remains flat at confirmation scale.
\item \textbf{Forgiveness / trust-repair experiment} (Section~\ref{sec:phenotypes}; Appendix~\ref{app:extended_results}): reengagement latency, post-return selection, payoff recovery, and whether high-$\alpha$ controls reengagement speed while cautious priors limit the trust-repair ceiling.
\item \textbf{H1 locality-probe confirmation}: replace focal-switch surprise--payoff correlation and aggregate payoff values with final 30-seed confirmation if they differ from current exploratory diagnostic values.
\item \textbf{H5 betrayal confirmation}: replace abrupt-betrayal payoff, entropy, and joint-accuracy values with final 30-seed H5 confirmation values when available.
\end{itemize}
